\documentclass[10pt,a4paper]{amsart}
\usepackage{amsmath,amssymb,amsthm,graphicx}
% Retained preamble for the Stacks Project chapter snapshot in this directory.
% The chapter's own preamble.tex is a copy of the shared Stacks Project preamble; the block
% below is sliced MECHANICALLY, line for line, from it. Nothing was added, removed or
% redefined; the wrapper class, the xr-hyper/hyperref wiring and the externaldocument list
% are omitted because the assembler supplies the document class.
%
% lines taken from preamble.tex: 3, 5, 9, 11, 12, 156, 157, 158, ... (51 lines total)
%
% package coverage restored from the chapter's own preamble (2026-10-13): xy, xy-2cell

\usepackage[all]{xy}
\xyoption{2cell}
\UseAllTwocells

\newenvironment{reference}{\comment}{\endcomment}
\newenvironment{slogan}{\comment}{\endcomment}
\newenvironment{history}{\comment}{\endcomment}
\theoremstyle{plain}
\newtheorem{theorem}[subsection]{Theorem}
\newtheorem{proposition}[subsection]{Proposition}
\newtheorem{lemma}[subsection]{Lemma}
\theoremstyle{definition}
\newtheorem{definition}[subsection]{Definition}
\newtheorem{example}[subsection]{Example}
\newtheorem{exercise}[subsection]{Exercise}
\newtheorem{situation}[subsection]{Situation}
\theoremstyle{remark}
\newtheorem{remark}[subsection]{Remark}
\newtheorem{remarks}[subsection]{Remarks}
\numberwithin{equation}{subsection}
\def\lim{\mathop{\mathrm{lim}}\nolimits}
\def\colim{\mathop{\mathrm{colim}}\nolimits}
\def\Spec{\mathop{\mathrm{Spec}}}
\def\Hom{\mathop{\mathrm{Hom}}\nolimits}
\def\Ext{\mathop{\mathrm{Ext}}\nolimits}
\def\SheafHom{\mathop{\mathcal{H}\!\mathit{om}}\nolimits}
\def\SheafExt{\mathop{\mathcal{E}\!\mathit{xt}}\nolimits}
\def\Sch{\mathit{Sch}}
\def\Mor{\mathop{\mathrm{Mor}}\nolimits}
\def\Ob{\mathop{\mathrm{Ob}}\nolimits}
\def\Sh{\mathop{\mathit{Sh}}\nolimits}
\def\NL{\mathop{N\!L}\nolimits}
\def\CH{\mathop{\mathrm{CH}}\nolimits}
\def\proetale{{pro\text{-}\acute{e}tale}}
\def\etale{{\acute{e}tale}}
\def\QCoh{\mathit{QCoh}}
\def\Ker{\mathop{\mathrm{Ker}}}
\def\Im{\mathop{\mathrm{Im}}}
\def\Coker{\mathop{\mathrm{Coker}}}
\def\Coim{\mathop{\mathrm{Coim}}}
\DeclareMathSymbol{\boxtimes}{\mathbin}{AMSa}{"02}
\def\QCohstack{\mathcal{QC}\!\mathit{oh}}
\def\Cohstack{\mathcal{C}\!\mathit{oh}}
\def\Spacesstack{\mathcal{S}\!\mathit{paces}}
\def\Quotfunctor{\mathrm{Quot}}
\def\Hilbfunctor{\mathrm{Hilb}}
\def\Curvesstack{\mathcal{C}\!\mathit{urves}}
\def\Polarizedstack{\mathcal{P}\!\mathit{olarized}}
\def\Complexesstack{\mathcal{C}\!\mathit{omplexes}}
\def\Pic{\mathop{\mathrm{Pic}}\nolimits}
\def\Picardstack{\mathcal{P}\!\mathit{ic}}
\def\Picardfunctor{\mathrm{Pic}}
\def\Deformationcategory{\mathcal{D}\!\mathit{ef}}
\usepackage{cleveref}
\crefname{originallemma}{original lemma}{original lemmas}
\Crefname{originallemma}{Original lemma}{Original lemmas}
\makeatletter
\newlabel{more-algebra-lemma-valuation-ring-torsion-free-flat}{{0539}{source}{}{sourcepointer-1}{}}
\newlabel{more-algebra-lemma-valuation-ring-torsion-free-flat@cref}{{[originallemma][0][]0539}{[1][0][]source}{}{}{}}
\makeatother
\crefname{originallemma}{original lemma}{original lemmas}
\Crefname{originallemma}{Original lemma}{Original lemmas}
\makeatletter
\newlabel{spaces-flat-theorem-finite-type-flat}{{0DLR}{source}{}{sourcepointer-2}{}}
\newlabel{spaces-flat-theorem-finite-type-flat@cref}{{[originallemma][0][]0DLR}{[1][0][]source}{}{}{}}
\makeatother
\crefname{originallemma}{original lemma}{original lemmas}
\Crefname{originallemma}{Original lemma}{Original lemmas}
\makeatletter
\newlabel{spaces-limits-lemma-directed-colimit-finite-type}{{0829}{source}{}{sourcepointer-3}{}}
\newlabel{spaces-limits-lemma-directed-colimit-finite-type@cref}{{[originallemma][0][]0829}{[1][0][]source}{}{}{}}
\makeatother
\crefname{originallemma}{original lemma}{original lemmas}
\Crefname{originallemma}{Original lemma}{Original lemmas}
\makeatletter
\newlabel{stacks-morphisms-definition-existence}{{0CLK}{source}{}{sourcepointer-4}{}}
\newlabel{stacks-morphisms-definition-existence@cref}{{[originallemma][0][]0CLK}{[1][0][]source}{}{}{}}
\makeatother
\title{Valuative existence for the Coh stack over a proper base}
\author{The Stacks Project Authors}
\date{stacks-project:moduli.tex:local-snapshot-20261012; GNU FDL 1.2 or later}
\begin{document}
\maketitle
\noindent\textit{Editorial note.} Statement span is lines 373-379 of sources/stacks-moduli/pinned-original-source.tex; proof span is lines 381-422, the complete author proof (40 body lines).\par\medskip
\noindent\textit{Editorial note.} Navigation stubs: 4 cited-but-not-included labels are supplied as non-load-bearing pointers with their pinned tags where the snapshot tag index has them.\par\medskip
\noindent\textit{Editorial note.} extra\_packages is empty on purpose: the retained preamble already loads xy with options and a second declaration clashes in the assembler wrapper.\par\medskip
\noindent\textit{Editorial note.} Text is retained verbatim, including the author's own typos and placeholders.\par\medskip
\noindent\textit{Editorial note.} Queue index 8 of queue\_moduli.json (pb=1945, stmt 373-379, proof 381-422).\par\medskip
\noindent\textit{Editorial note.} stacks-moduli wave-1: fresh adjudication accept (2026-10-14); pool per Chief order 08:5x (dpa superseded to lane-s9); moduli.tex 53 labels/48 with proof judged 25 accept / 23 reject / 5 pb=0 env-dup rows excluded; spans cite-free and reference/slogan/history-free so no patchenv needed.\par\medskip
\noindent\textit{Editorial note.} Author's '(small detail omitted)' covers the routine adjunction identification j\textasciicircum{}*j\_*F\_K = F\_K; the argument chain is complete.\par\medskip
\noindent\textit{Allowed standard prerequisites (editorial).}\begin{itemize}
\item Explicit prerequisite boundary. The cited inputs are prerequisites; the argument in the retained span is not replaced by any cited result.
\end{itemize}\medskip
\ifdefined\hypertarget\hypertarget{sourcepointer-1}{}\fi
\noindent\textit{Source reference.} Original lemma 0539: sources/stacks-moduli/pinned-original-source.tex, source lines 422--422. Cited by the retained statement or proof but not part of the collected spans; supplied only as a navigation pointer. The argument does not depend on its content.\par\medskip
\ifdefined\hypertarget\hypertarget{sourcepointer-2}{}\fi
\noindent\textit{Source reference.} Original lemma 0DLR: sources/stacks-moduli/pinned-original-source.tex, source lines 422--422. Cited by the retained statement or proof but not part of the collected spans; supplied only as a navigation pointer. The argument does not depend on its content.\par\medskip
\ifdefined\hypertarget\hypertarget{sourcepointer-3}{}\fi
\noindent\textit{Source reference.} Original lemma 0829: sources/stacks-moduli/pinned-original-source.tex, source lines 422--422. Cited by the retained statement or proof but not part of the collected spans; supplied only as a navigation pointer. The argument does not depend on its content.\par\medskip
\ifdefined\hypertarget\hypertarget{sourcepointer-4}{}\fi
\noindent\textit{Source reference.} Original lemma 0CLK: sources/stacks-moduli/pinned-original-source.tex, source lines 422--422. Cited by the retained statement or proof but not part of the collected spans; supplied only as a navigation pointer. The argument does not depend on its content.\par\medskip
\begin{lemma}
\label{lemma-coherent-existence-part}
Assume $X \to B$ is proper as well as of finite presentation.
Then $\Cohstack_{X/B} \to B$ satisfies the existence part
of the valuative criterion (Morphisms of Stacks, Definition
\ref{stacks-morphisms-definition-existence}).
\end{lemma}

\begin{proof}
Taking base change, this immediately reduces to the following
problem: given a valuation ring $R$ with fraction field $K$ and
an algebraic space $X$ proper over $R$ and a coherent
$\mathcal{O}_{X_K}$-module $\mathcal{F}_K$, show there exists
a finitely presented $\mathcal{O}_X$-module $\mathcal{F}$
flat over $R$ whose generic fibre is $\mathcal{F}_K$.
Observe that by Flatness on Spaces, Theorem
\ref{spaces-flat-theorem-finite-type-flat}
any finite type quasi-coherent $\mathcal{O}_X$-module
$\mathcal{F}$ flat over $R$ is of finite presentation.
Denote $j : X_K \to X$ the embedding of the generic fibre.
As a base change of the affine morphism $\Spec(K) \to \Spec(R)$
the morphism $j$ is affine. Thus $j_*\mathcal{F}_K$ is
quasi-coherent. Write
$$
j_*\mathcal{F}_K = \colim \mathcal{F}_i
$$
as a filtered colimit of its finite type quasi-coherent
$\mathcal{O}_X$-submodules, see
Limits of Spaces, Lemma \ref{spaces-limits-lemma-directed-colimit-finite-type}.
Since $j_*\mathcal{F}_K$ is a sheaf of $K$-vector spaces over $X$,
it is flat over $\Spec(R)$. Thus each $\mathcal{F}_i$ is flat
over $R$ as flatness over a valuation ring is the same as being
torsion free
(More on Algebra, Lemma
\ref{more-algebra-lemma-valuation-ring-torsion-free-flat})
and torsion freeness is inherited by submodules.
Finally, we have to show that the map
$j^*\mathcal{F}_i \to \mathcal{F}_K$
is an isomorphism for some $i$.
Since $j^*j_*\mathcal{F}_K = \mathcal{F}_K$ (small detail omitted)
and since $j^*$ is exact, we see that $j^*\mathcal{F}_i \to \mathcal{F}_K$
is injective for all $i$.
Since $j^*$ commutes with colimits, we have
$\mathcal{F}_K = j^*j_*\mathcal{F}_K = \colim j^*\mathcal{F}_i$.
Since $\mathcal{F}_K$ is coherent (i.e., finitely presented),
there is an $i$ such that $j^*\mathcal{F}_i$ contains all the
(finitely many) generators over an affine \'etale cover of $X$.
Thus we get surjectivity of $j^*\mathcal{F}_i \to \mathcal{F}_K$
for $i$ large enough.
\end{proof}

\end{document}
